\section{Related Work}

The VOC-ALS dataset was introduced by Dubbioso et al. and contains 1,224 smartphone recordings from 153 participants: 102 people with ALS and 51 healthy controls \cite{dubbioso2024vocals}. Each participant completed five sustained vowel tasks and three syllable-repetition tasks. The dataset also includes clinical information and several basic acoustic measurements, including fundamental frequency, jitter, shimmer, and harmonics-to-noise ratio.

Previous studies have shown that speech features can capture useful information about ALS-related speech changes. Simmatis et al. used automatically extracted acoustic features to distinguish people with ALS from healthy controls and to separate different levels of bulbar impairment \cite{simmatis2024acoustic}. More recent work has also compared global and vowel-level acoustic features for dysarthria assessment in motor neuron disease \cite{attia2025global}.

These studies show that speech contains measurable information related to motor impairment. My project focuses on a different question: how much speech is actually needed. Instead of only asking which model performs best when using all available tasks, I will compare every possible subset of the eight VOC-ALS tasks and look at the tradeoff between model performance and the number of tasks required.

